\slajdid{08-s049}
\begin{frame}[label=08-s049]
\begin{columns}[c,onlytextwidth]\column{.53\textwidth}\centering\input{slike/tikz/s049_eksponencijalni_odziv.tex}
\column{.44\textwidth}\gusttekst
Gledajući izraz za izlazni napon, vidimo da vrednost teži ka konačnoj vrednosti U, i dostiže je tek u tački $t=\infty$, kada se i završava prelazni proces. Za praktične slučajeve a i za crtanje, ovo nam baš i nije zgodno.\par\bigskip
\textbf{Zbog toga se u praksi smatra da se prelazni proces završava posle vremena \alert{$5\tau$}, kada se dostiže 99.3\% promene.}
\end{columns}
\end{frame}
